\documentclass[oneside]{book}
%\documentclass[onecolumn]{emulateapj-rtx4}
\usepackage{natbib}
\usepackage{color}
\usepackage{apjfonts}
\usepackage{graphicx}
\usepackage{amssymb}
\usepackage{amsbsy}
\usepackage{amsmath}
\usepackage{mathrsfs}
\usepackage{geometry}
\geometry{left=1.5cm, right=1.5cm, top=2.5cm, bottom=2.0cm}

\def\rd{{\rm d}}
\def\mnras{{MNRAS}}
\def\apj{{ApJ}}
\def\apjl{{ApJ}}
\def\aap{{A\&A}}
\begin{document}

\title{Structure and Emergent Spectrum of Slim Accretion Disks With Hot Coronae}
\author{Yan-Rong Li\\
liyanrong@mail.ihep.ac.cn\\
Institute of High Energy Physics}
\date{}
\maketitle

\chapter{Basic Equations}
To mimic the general relativistic effects, we adopt the pseudo-Newtonian potential
\begin{equation}
\Phi(R)=-\frac{GM_\bullet}{R-2R_{\rm g}},
\end{equation}
where $R_{\rm g}=GM_\bullet/c^2$ is the gravitational radius. The inner edges of slim disks locate at 
$R_{\rm in}=6R_{\rm g}$, where the usual ``torque-free'' condition is enforced. The vertical structures of slim disks
and hot coronae are in hydrodynamical equilibrium and the vertical temperature profile is isothermal.
The geometry of the coronae is assumed to be parallel-plane above the disk.
\subsection{Slim Accretion Disks}
The governing equations for mass, angular momentum, radial velocity, and energy are as follows
(\citealt{Wang2003}).
\begin{itemize}
 \item Mass
 %
\begin{equation}
\dot M=-2\pi R V_{\rm R}\Sigma,
\label{equ_disk_mass}
\end{equation}
%
where $\dot M$ is the mass accretion rate, $V_R$ is the radial velocity, and $\Sigma$ is the surface density.
\item Angular momentum
%
\begin{equation}
\dot M(\ell-\ell_{\rm in})=-2\pi R^2 T_{r\phi},
\label{equ_disk_am}
\end{equation}
where $\ell$ ($\ell_{\rm in}$) is the specific angular momentum (at the inner edge) and $T_{r\phi}$
is the height-integrated shear tensor.
\item Radial velocity
%
\begin{equation}
V_R\frac{\rd V_R}{\rd R}+\frac{1}{\Sigma}\frac{\rd W}{\rd R}
=\frac{\ell^2-\ell_{\rm K}^2}{R^3}-\frac{W}{\Sigma}\frac{\rd\ln\Omega_{\rm K}}{\rd R},
\label{equ_disk_vel}
\end{equation}
where $W$ is the height-integrated pressure, $\ell_{\rm K}=R^2\Omega_{\rm K}$, and $\Omega_{\rm K}$
is the Kerplerian angular velocity.
\item Energy
%
\begin{equation}
Q_{\rm vis}^+=Q_{\rm rad}^-+Q_{\rm adv}^-+Q_{\rm cor}^+ - \varepsilon(1-a_r)Q_{\rm cor}^+,
\label{equ_disk_energy}
\end{equation}
where $\varepsilon$ is the fraction of radiation that are intercepted by the disks, $a_r$ is the reflection albedo,
$Q_{\rm vis}^+$ is the viscous heating, $Q_{\rm rad}^-$ is the radiative cooling, 
$Q_{\rm adv}^-$ is the advective cooling, and $Q_{\rm cor}^+$ is the energy 
dissipated into the hot coronae. 
\end{itemize}
Basically, we write down the expresses for the related terms in the energy equation as
\begin{gather}
Q_{\rm vis}^+=RT_{r\phi}\frac{\rd\Omega}{\rd r},\\
Q_{\rm rad}^-=\frac{8acT_0^4}{3\kappa\rho_0H_d},\\
Q_{\rm adv}^-=\frac{\dot M}{4\pi R^2}\frac{W}{\Sigma}\xi,\\
Q_{\rm cor}^+=\frac{B^2}{8\pi}v_p,
\end{gather}
where 
\begin{equation}
\kappa=0.40+0.64\times10^{23}\rho_0 T_0^{-7/2} g^{-1}{\rm cm}^2,
\end{equation}
and 
\begin{equation}
\xi=(4-3\beta)\frac{\rd\ln\Sigma}{\rd\ln R}-\left(12-\frac{21}{2}\beta\right)\frac{\rd\ln T}{\rd\ln R}-
(4-3\beta)\frac{\rd\ln H_d}{\rd\ln R},
\end{equation}
with $\beta=P_g/(P_g+P_r)$.
The viscous stress tensor $T_{r\phi}=-\alpha_{\rm eff}W$ and $W=\int P_t d z$ is the height-integrated pressure.
The total pressure is
\begin{equation}
 P_t=P_g+P_r=\frac{\rho_0 k T_0}{\bar\mu m_H} + \frac{1}{3}aT_0^4.
\end{equation}
The effective coefficient $\alpha_{\rm eff}=\alpha(P_g/P_t)^\mu$. The magnetic pressure is (\citealt{Merloni2003})
\begin{equation}
P_m=\alpha_0\sqrt{P_gP_t},
\end{equation}
where $\alpha_0$ is order of unity. The velocity of the magnetic flux transported vertially in the disk is in proportional
to the Alfven velocity $v_p=bv_A$ with $v_A=B/\sqrt{4\pi\rho}$ and $b\approx1$.
The fraction of the viscous heating dissipated to the coronae is
\begin{equation}
f(R)=\frac{Q_{\rm cor}^+}{Q_{\rm vis}^+}.
\end{equation}
The effective temperature of the disk is
\begin{equation}
Q_{\rm rad}^-=2F_d^-=2\sigma T_{\rm eff}^4.
\end{equation}

(There are {\color{blue} five} variables to be determined: 
{\color{blue}$\Sigma$, $V_R$, $W$, $\Omega$ and $H_d$}.)


\subsection{Hot Coronae}
For hot coronae, the energy equations for ions and electrons are as follows.
\begin{itemize}
 \item Ions
 \begin{gather}
Q_{\rm cor}^+-\lambda Q_{\rm cor}^+=Q_{ie}^-.
\label{equ_cor_ions}
\end{gather}
 \item Electrons
 \begin{equation}
  Q_{ie}^-+\lambda Q_{\rm cor}^+=2F_c^-.
  \label{equ_cor_electrons}
 \end{equation}
\end{itemize}
where $\lambda$ is the fraction of energy that directly heats the electrons, $Q_{ie}^-$ is 
energy transfer between ions and electrons via the Coulomb interactions.
and $F_c^-$ is the cooling rate of the coronae including the synchrotron, bremstrahlung emssions and
the Compton scattering.
The scale height of the hot coronae is
\begin{equation}
P_c=\rho_c\Omega_K^2H_c^2,
\label{equ_cor_height}
\end{equation}
where $P_c=P_{c, g}+P_{c, m}$ is the total pressure of the coronae.  The magnetic pressure is assumed to a fraction
of the totoal pressure $P_{c, m}=B_c^2/8\pi=\beta_c P_c$, and 
\begin{equation}
P_{c, g}=\frac{\rho_ckT_i}{\mu_i m_H}+\frac{\rho_ckT_e}{\mu_e m_H}.
\end{equation}

(There are {\color{blue} four} variables to be determined: {\color{blue}$T_i$, $T_e$, $n$, and $H_c$}.)

\section{Disk-Coronae Transition}
\subsection{Model I: Pressure Balance}
There is pressure balance at the base of the coronae. Therefore, the height of the disk satisfies
\begin{equation}
\Sigma H_d^2\Omega_{\rm K}^2=W+P_mH_d-P_{\rm base}H_d.
\label{equ_disk_height}
\end{equation}
where $P_{\rm base}$ is the pressure at the base of the coronae.
Following \cite{Zycki1995}, the tansition from the disk to the coronae can be qualitatively determined by 
requirement that the ionization parameter
\begin{equation}
\Xi\equiv \frac{F_d}{cP_{\rm g}}=\Xi_{\rm min},
\end{equation}
where $F_d$ is the radiation flux from the disk and {\color{blue}$\Xi_{\rm min}=0.05$}.
Similarly, at the base of the coronae, there should be
\begin{equation}
P_c=P_{\rm base}.
\label{equ_cor_base}
\end{equation}

The equation sets for the disk and the coronae are now closed.

\subsection{Model II: Evaporation and Thermal Conduction}
The height of the disk is given by neglecting the pressure of the coronae
\begin{equation}
\Sigma H_d^2\Omega_{\rm K}^2=W.
\end{equation}
For the coronae, at its interface, there should be energy balance between the 
evaporation and thermal conduction (\citealt{Liu2002})
\begin{equation}
\frac{\kappa_0 T^{7/2}}{l}=\frac{\gamma}{\gamma-1}nkT\left(\frac{kT}{m_H}\right)^{1/2},
\end{equation}
where {\color{blue}$l\approx10R_S$} is the characteristic length of the magnetic loops in the coronae.

The equation sets for the disk and the coronae are now closed.

\section{Radiation Meachnism}
The energy transfer between the ions and electrons $Q_{\rm ie}^-$ is given by (\citealt{Narayan1995})
\begin{equation}
Q_{\rm ie}^-=H_c\times1.25\frac{3}{2}\frac{m_e}{m_p}
n_en_i\sigma_T c\frac{kT_i-kT_e}{K_2(1/\theta_e)K_2(1/\theta_i)}\ln\Lambda
\left[\frac{2(\theta_e+\theta_i)^2+1}{\theta_i+\theta_e}K_1\left(\frac{\theta_e+\theta_i}{\theta_e\theta_i}\right)
+2K_0\left(\frac{\theta_e+\theta_i}{\theta_e\theta_i}\right)\right]~~{\rm erg~s^{-1}~cm^{-2}}.
\end{equation}
The synchrotron emission is calculated by $F_{\rm syn}^-=H_c \chi_{syn}$
\begin{eqnarray}
\chi_{syn}=4.43\times10^{-30}\frac{4\pi n_e\nu}{K_2(1/\theta_e)}
I\left(\frac{4\pi m_ec\nu}{3eB\theta_e^2}\right),
\end{eqnarray}
where
\begin{eqnarray}
I(x)=\frac{4.0505}{x^{1/6}}\left(1+\frac{0.4}{x^{1/4}}+\frac{0.5316}{x^{1/2}}\right)
\exp(-1.8899x^{1/3}).
\end{eqnarray}
The bremstralung emission is calculated  by $F_{\rm br}^-=H_c \chi_{br}$ with
\begin{eqnarray}
\chi_{br}=q_{br}^{-}\bar G\exp\left(\frac{h\nu}{kT_e}\right),
\end{eqnarray}
where $\bar G$ is the Gaunt factor, $q_{br}^{-}$ is the bremsstrahlung emission per unit volume
, which reads
\begin{eqnarray}
q_{br}^{-}=1.48\times10^{-22}n_e^2F(\theta),
\end{eqnarray}
and 
\begin{eqnarray}
F(\theta)&&=4\left(\frac{2\theta_e}{\pi^3}\right)^{1/2}(1+1.781\theta_e^{1.34})
+1.73\theta_e^{3/2}(1+1.1\theta_e+\theta_e^2-1.25\theta_e^{5/2})\qquad\theta_e<1,\nonumber\\
&&=\left(\frac{9\theta_e}{2\pi}\right)[\ln(0.48+1.123\theta_e)+1.5]
+2.30\theta_e(\ln1.123\theta_e+1.28)\qquad\theta_e>1.
\end{eqnarray}
The Gaunt factor is given by (\citealt{Rybicki1979})
\begin{equation}
\bar G=\frac{h}{kT_e}\left(\frac{3}{\pi}\frac{kT_e}{h\nu}\right)^{1/2}\qquad \frac{kT_e}{h\nu}<1,
\end{equation}
\begin{equation}
\bar G=\frac{h}{kT_e}\frac{\sqrt{3}}{\pi}\ln\left(\frac{4}{\xi}\frac{kT_e}{h\nu}\right)^{1/2}
\qquad \frac{kT_e}{h\nu}>1.
\end{equation}
We treat the Compton scattering as follows.
The energy enhancement factor $\eta$ is
\begin{equation}
 \eta=\exp[s(A-1)][1-P(j_m+1,As)]+\eta_{max}P(j_m+1,s),
\end{equation}
where $P$ is the incomplete gamma function and
\begin{equation}
 A=1+4\theta_e+16\theta_e^2,~~~~~s=\tau_{es}+\tau_{es}^2,
\end{equation}
and
\begin{equation}
 \eta_{max}=\frac{3kT_e}{h\nu},~~~~~~~~~~j_m=\frac{\ln\eta_{max}}{\ln A},
~~~~~~~~~~\tau_{es}=n_e\sigma_TH_c.
\end{equation}
With the energy enhancement factor $\eta$, the local radiative cooling rate $Q_{rad}^{-}$ is given by
\begin{equation}
F_{c}^{-}=2\int d\nu \left[\eta (F_{\rm br, \nu} + F_{\rm syn, \nu}) + (\eta-1) F_{d, \nu}\right],
\end{equation}
where $F_\nu$ includes the local synchrotron and bremstrahlung emission and the flux from the accretion disk.
\section{Emergent Spectrum: Calculations}
The emergent spectrum from slim disks and hot coronae is calculated as
\begin{equation}
F_\nu=F_{d,\nu} +(1-\varepsilon) \left(F_{\rm syn, \nu} + F_{\rm br, \nu} + F_{Comp, \nu}\right),
\end{equation}
where $\varepsilon$ is the fraction of emissions that are intercepted by the disks. 
The spectra from Compton scatterings in the coronae is treated according to \cite{Coppi1990}.

\chapter{Derive the Equations}
The basic equations are constructed in the cyclindrical coordinates $(R, \phi, z)$.
The mass conservation gives
\begin{equation}
\frac{\partial (R\rho v_r)}{R\partial R}+\frac{\partial (\rho v_z)}{\partial z}=0.
\end{equation}
The radial component of  the momentum equation reads
\begin{equation}
\frac{\partial (R\rho v_r^2)}{R\partial R}+\frac{\partial(\rho v_r v_z)}{\partial z}-\frac{\rho v_\phi^2}{R}
=-\frac{\partial P}{\partial R}-\rho\frac{\partial\Phi}{\partial R}.
\end{equation}
The azimuthal component of the momentum equation is
\begin{equation}
\frac{\partial(R^2\rho v_rv_\phi)}{R\partial R}+\frac{\partial (\rho v_\phi v_z)}{\partial z}
=\frac{\partial (R^2 \tau_{r\phi})}{R\partial R}.
\end{equation}
The energy equation is 
\begin{equation}
 \rho T\left(v_r\frac{\partial s}{\partial R}+v_z\frac{\partial s}{\partial z}\right)
 =t_{r\phi} r\frac{\rd\Omega}{\rd R}-q_{\rm rad}^-.
\end{equation}
The hydrostatic balance in the vertical direction is
\begin{equation}
\frac{\partial P}{\partial z}=-\rho\frac{\partial \Phi}{\partial z}\approx-\rho\Omega_{\rm K}^2z.
\end{equation}
If we assume the disk is isothermal in the vertical direction, the solution is 
\begin{equation}
\rho(R, z)=\rho_0(R)\exp\left(-\frac{z^2}{2H^2}\right),
\end{equation}
with $H=c_s/\Omega_{\rm K}$.

Define following variables
\begin{equation}
\Sigma=\int\rho\rd z, ~~~W=\int P\rd z,
\end{equation}
we get the mass accretion rate
\begin{equation}
 \dot M=-2\pi R\Sigma v_r={\rm const.}
\end{equation}
The radial component of the momentum equation is 
\begin{equation}
v_r\frac{\rd v_r}{\rd R}+\frac{1}{\Sigma}\frac{\rd W}{\rd R}=
\frac{\ell^2-\ell_{\rm K}^2}{R^3}-\frac{W}{\Sigma}\frac{\rd\ln\Omega_{\rm K}}{\rd R}.
\end{equation}
The azimuthal component of the momentum equation is
\begin{equation}
\dot M(\ell-\ell_{\rm in})=-2\pi R^2 T_{r\phi}.
\end{equation}

We present the thermal dynamical relations (\citealt{Kato2008})
\begin{equation}
\rho T \rd s=\frac{1}{\Gamma_3-1}\left(\rd P -\Gamma_1\frac{P}{\rho}\rd \rho\right),
\end{equation}
where 
\begin{equation}
\Gamma_1=\beta + \frac{(\gamma-1)(4-3\beta)^2}{\beta + 12(\gamma-1)(1-\beta)},~~~~
\Gamma_3=1+\frac{(4-3\beta)(\gamma-1)}{\beta + 12(\gamma-1)(1-\beta)}.
\end{equation}
As a result, the left-hand side of the energy equation can be recast as
\begin{equation}
 \rho T\left(v_r\frac{\partial s}{\partial R}+v_z\frac{\partial s}{\partial z}\right)
 =\frac{1}{\Gamma_3-1}\left[v_r\frac{\partial P}{\partial R} + v_z\frac{\partial P}{\partial z}
 -\Gamma_1\frac{P}{\rho}\left(v_r\frac{\partial \rho}{\partial R} + v_z\frac{\partial \rho}{\partial z}\right)\right].
\end{equation}
With help of the continuity equation, 
\begin{equation}
\frac{P}{\rho}\left(v_r\frac{\partial \rho}{\partial R} + v_z\frac{\partial \rho}{\partial z}\right)
=Pv_r+P\frac{\partial v_r}{\partial R} + P \frac{\partial v_z}{\partial z}.
\end{equation}
This simplifies the above equation into
\begin{equation}
 \rho T\left(v_r\frac{\partial s}{\partial R}+v_z\frac{\partial s}{\partial z}\right)
 =\frac{\Gamma_1}{\Gamma_3-1}\left[\frac{\partial (Rv_rP)}{R\partial R} +
 \frac{\partial (v_zP)}{\partial z}\right]
 -\frac{\Gamma_1-1}{\Gamma_3-1}\left[v_r\frac{\partial P}{\partial R}+v_z\frac{\partial P}{\partial z}\right].
\end{equation}
After the vertial integrating, we have 
\begin{gather}
\frac{\Gamma_1}{\Gamma_3-1}v_r\Sigma\frac{\rd (W/\Sigma)}{\rd R}
-\frac{\Gamma_1-1}{\Gamma_3-1}\left[v_r\frac{\rd W}{\rd R}-v_rW\frac{\rd\ln H}{\rd R}\right]\\
=-\frac{1}{\Gamma_3-1}\frac{\dot M}{2\pi R}\frac{W}{\Sigma}
\left[\frac{\rd \ln W}{\rd R} -\Gamma_1\frac{\rd \ln \Sigma}{\rd R}+(\Gamma_1-1)\frac{\rd\ln H}{\rd R}\right]\\
=-\frac{1}{\Gamma_3-1}\frac{\dot M}{2\pi R}\frac{W}{\Sigma}
\left[\frac{\Gamma_1+1}{2}\frac{\rd \ln W}{\rd R}-
\frac{3\Gamma_1-1}{2}\frac{\rd \ln \Sigma}{\rd R}
-(\Gamma_1-1)\frac{\rd\ln\Omega_{\rm K}}{\rd R}\right].
\end{gather}
The energy equation is rewritten into
\begin{equation}
-\frac{1}{\Gamma_3-1}\frac{\dot M}{2\pi R}\frac{W}{\Sigma}
\left[\frac{\Gamma_1+1}{2}\frac{\rd \ln W}{\rd R}-
\frac{3\Gamma_1-1}{2}\frac{\rd \ln \Sigma}{\rd R}
-(\Gamma_1-1)\frac{\rd\ln\Omega_{\rm K}}{\rd R}\right]
=RT_{r\phi}\frac{\rd\Omega}{\rd R}-Q_{\rm rad}^-.
\end{equation}
It can be further recasted as
\begin{equation}
\frac{\dot M}{\Gamma_3-1}\frac{W}{\Sigma}
\left[\frac{\Gamma_1+1}{2}\frac{\rd \ln W}{\rd R}-
\frac{3\Gamma_1-1}{2}\frac{\rd \ln \Sigma}{\rd R}
-(\Gamma_1-1)\frac{\rd\ln\Omega_{\rm K}}{\rd R}\right]
=\dot M (\ell-\ell_{\rm in})\frac{\rd\Omega}{\rd R}+2\pi R Q_{\rm rad}^-.
\end{equation}

In a summary, the equation set for the disk is 
\begin{gather}
 \dot M=-2\pi R\Sigma v_r={\rm const},\\
\dot M(\ell-\ell_{\rm in})=-2\pi R^2 T_{r\phi},\\
v_r\frac{\rd v_r}{\rd R}+\frac{1}{\Sigma}\frac{\rd W}{\rd R}=
\frac{\ell^2-\ell_{\rm K}^2}{R^3}-\frac{W}{\Sigma}\frac{\rd\ln\Omega_{\rm K}}{\rd R},\\
\frac{\dot M}{\Gamma_3-1}\frac{W}{\Sigma}
\left[\frac{\Gamma_1+1}{2}\frac{\rd \ln W}{\rd R}-
\frac{3\Gamma_1-1}{2}\frac{\rd \ln \Sigma}{\rd R}
-(\Gamma_1-1)\frac{\rd\ln\Omega_{\rm K}}{\rd R}\right]
=\dot M (\ell-\ell_{\rm in})\frac{\rd\Omega}{\rd R}+2\pi R Q_{\rm rad}^-.
\end{gather}

\chapter{Solving the Equations}
\section{Numerical Scheme}
We choose $\Sigma$ and $W$ as the independent variables. The shear tensor is
\begin{equation}
T_{r\phi}=-\alpha_{\rm eff}W, ~~~\alpha_{\rm eff}=\alpha(P_g/P_t)^\mu.
\end{equation}
The derivation of the scale height is 
\begin{equation}
\frac{\rd\ln H}{\rd R}=\frac{1}{2}\frac{\rd\ln W}{\rd R}-\frac{1}{2}\frac{\rd\ln\Sigma}{\rd R}
-\frac{\rd\ln\Omega_{\rm K}}{\rd R}.
\end{equation}
The derivation of the radial velocity is
\begin{equation}
v_r\frac{\rd v_r}{\rd R}=-\frac{\dot M^2}{4\pi^2 R^2\Sigma^2}\left[\frac{1}{R}+\frac{\rd\ln\Sigma}{\rd R}\right],
\end{equation}
and the derivation of the angular velocity is 
\begin{equation}
\frac{\rd\Omega}{\rd R}=-\frac{\ell_{\rm in}}{R^3}+\frac{2\pi}{\dot M}\frac{\rd}{\rd R}\left(\alpha_{\rm eff} W\right)
=-\frac{\ell_{\rm in}}{R^3}+\frac{2\pi\alpha_{\rm eff} W}{\dot M}\frac{\rd}{\rd R}\left[\mu\ln W_g +(1-\mu)\ln W\right].
\end{equation}
Using the equation of state, we have 
\begin{equation}
W=W_g+W_r=\frac{\Sigma k T_0}{\bar\mu m_H} + \frac{2H}{3}aT_0^4,
\end{equation}
and
\begin{equation}
\frac{\rd W}{\rd R}=W_g\frac{\rd \ln\Sigma}{\rd R} + W_g\frac{\rd\ln T_0}{\rd R}
+4W_g\frac{\rd\ln T_0}{\rd R} + W_r\frac{\rd \ln H}{\rd R}.
\end{equation}
Thie yields
\begin{gather}
\frac{\rd\ln T_0}{\rd R}=\frac{1}{4-3\beta}\left[\frac{\rd\ln W}{\rd R}
-\beta\frac{\rd\ln\Sigma}{\rd R} - (1-\beta)\frac{\rd \ln H}{\rd R}\right]\\
=\frac{1}{4-3\beta}\left[\frac{\beta +1}{2}\frac{\rd\ln W}{\rd R}
+\frac{1-3\beta}{2}\frac{\rd\ln\Sigma}{\rd R}+ (1-\beta)\frac{\rd \ln \Omega_{\rm K}}{\rd R}\right],
\end{gather}
and 
\begin{gather}
\frac{\rd \ln W_g}{\rd R}=\frac{1}{\beta}\frac{\rd\ln W}{\rd R} - \frac{1-\beta}{\beta}\frac{\rd\ln W_r}{\rd R}
=\frac{1}{\beta}\frac{\rd\ln W}{\rd R} - \frac{1-\beta}{\beta}\left(\frac{\rd\ln H}{\rd R}+4\frac{\rd\ln T_0}{\rd R}\right)\\
=\frac{1+\beta}{2(4-3\beta)}\frac{\rd\ln W}{\rd R}
+\frac{9(1-\beta)}{2(4-3\beta)}\frac{\rd\ln \Sigma}{\rd R}
+\frac{1-\beta}{(4-3\beta)}\frac{\rd\ln \Omega_{\rm K}}{\rd R}.
\end{gather}
With above equatuion, we have
\begin{gather}
\frac{\rd\Omega}{\rd R}=-\frac{\ell_{\rm in}}{R^3}+\frac{2\pi\alpha_{\rm eff} W}{\dot M}
\left\{\left[\frac{\mu(1+\beta)}{2(4-3\beta)}+(1-\mu)\right]\frac{\rd\ln W}{\rd R}
+\frac{9\mu(1-\beta)}{2(4-3\beta)}\frac{\rd\ln \Sigma}{\rd R}
+\frac{\mu(1-\beta)}{(4-3\beta)}\frac{\rd\ln \Omega_{\rm K}}{\rd R}
\right\}.
\end{gather}
The radiative cooling is
\begin{equation}
Q_{\rm rad}^-=\frac{8c(1-\beta)}{\kappa}\frac{W}{\Sigma H}.
\end{equation}

With above equations, we have the equation set for numerical solving
\begin{gather}
\frac{W}{\Sigma}\frac{\rd\ln W}{\rd R}-\frac{\dot M^2}{4\pi^2R^2\Sigma^2}\frac{\rd\ln\Sigma}{\rd R}
=\frac{\ell^2-\ell_{\rm K}^2}{R^3}-\frac{W}{\Sigma}\frac{\rd\ln\Omega_{\rm K}}{\rd R}+
\frac{\dot M^2}{4\pi^2R^3\Sigma^2},\\
\left\{ \frac{\Gamma_1+1}{2(\Gamma_3-1)}\frac{\dot MW}{\Sigma}
 -2\pi\alpha_{\rm eff}(\ell-\ell_{\rm in})W\left[\frac{\mu(1+\beta)}{2(4-3\beta)}+(1-\mu)\right]\right\}
\frac{\rd\ln W}{\rd R}\nonumber\\
-\left\{ \frac{3\Gamma_1-1}{2(\Gamma_3-1)}\frac{\dot MW}{\Sigma}
+2\pi\alpha_{\rm eff}(\ell-\ell_{\rm in})W\frac{9\mu(1-\beta)}{2(4-3\beta)}\right\}
\frac{\rd\ln \Sigma}{\rd R}\nonumber\\
=-\dot M(\ell-\ell_{\rm in})\frac{\ell_{\rm in}}{R^3}+2\pi R Q_{\rm rad}^-
+\left\{\frac{\Gamma_1-1}{\Gamma_3-1}\frac{\dot MW}{\Sigma}
+2\pi\alpha_{\rm eff}(\ell-\ell_{\rm in})W\frac{\mu(1-\beta)}{(4-3\beta)}\right\}
\frac{\rd\ln \Omega_{\rm K}}{\rd R}
\end{gather}
The solutions are
\begin{gather}
\frac{\rd\ln W}{\rd R}=\frac{a_{22}c_1-a_{12}c_2}{a_{11}a_{22}-a_{21}a_{12}},\\
\frac{\rd\ln \Sigma}{\rd R}=\frac{a_{11}c_2-a_{21}c_1}{a_{11}a_{22}-a_{21}a_{12}}.
\end{gather}
Once $W$ and $\Sigma$ are obtained, other variabes are calculated as
\begin{gather}
\Omega=\frac{\ell_{\rm in}}{R^2}+\frac{2\pi\alpha_{\rm eff}W}{\dot M}, ~~~~
H\approx\frac{1}{\Omega_{\rm K}}\sqrt{\frac{W}{\Sigma}},~~~~
W=\frac{\Sigma k T_0}{\bar\mu m_H} + \frac{2H}{3}aT_0^4
\end{gather}

The transonic point is located where the denominator $\mathscr{D}=a_{11}a_{22}-a_{21}a_{12}=0$. 
The denominator has a form of 
\begin{eqnarray}
\frac{\Sigma}{W\dot M}\mathscr{D}&=&-\frac{W}{\Sigma}
\left[ \frac{3\Gamma_1-1}{2(\Gamma_3-1)} 
+\frac{4\pi^2\alpha_{\rm eff}^2 R^2 W\Sigma}{\dot M^2} \frac{9\mu(1-\beta)}{2(4-3\beta)}\right]\\
&&+\frac{\dot M^2}{4\pi^2R^2\Sigma^2}\left\{\frac{\Gamma_1+1}{2(\Gamma_3-1)}
-\frac{4\pi^2\alpha_{\rm eff}^2 R^2 W\Sigma}{\dot M^2}\left[\frac{\mu(1+\beta)}{2(4-3\beta)}+(1-\mu)\right]\right\}\\
&=&-c_s^2\left(b_1 + b_2\frac{c_s^2}{v_r^2}\right)
+v_r^2\left(b_3-b_4\frac{c_s^2}{v_r^2}\right)\\
&=&-v_r^2\left[b_2\left(\frac{c_s^2}{v_r^2}\right)^2 + (b_1+b_4)\frac{c_s^2}{v_r^2} - b_3\right].
\end{eqnarray}
The solution is 
\begin{equation}
\frac{c_s^2}{v_r^2}=\frac{-(b_1+b_4)+\sqrt{(b_1+b_4)^2+4b_2b_3}}{2b_2}.
\end{equation}
where 
\begin{gather}
b_1=\frac{3\Gamma_1-1}{2(\Gamma_3-1)},~~~~b_2= \alpha_{\rm eff}^2\frac{9\mu(1-\beta)}{2(4-3\beta)}\\
b_3=\frac{\Gamma_1+1}{2(\Gamma_3-1)},~~~~b_4= \alpha_{\rm eff}^2\left[\frac{\mu(1+\beta)}{2(4-3\beta)}+(1-\mu)\right]
\end{gather}

In the calcuations, we adopt the dimensional units
\begin{gather}
R: R_g,~~~\ell: R_gc,~~~\Omega:\frac{c}{R_g},~~~\dot M:\dot M_{\rm Edd},~~~\Sigma:\frac{\dot M_{\rm Edd}}{cR_g}\\
W: \frac{\dot M_{\rm Edd} c}{R_g}, ~~~Q_{\rm rad}^-=\frac{\dot M_{\rm Edd} c^2}{R_g^2}
\end{gather}

The boundary conditions are set according to \cite{Wang1999}
\begin{equation}
\Omega=\frac{\Omega_{\rm K}}{\sqrt{5}},~~~T=1.37\times10^8\left(\frac{\dot m}{16m}\right)^{1/4}
\left(\frac{10R_g}{R}\right)^{5/8}\left(\frac{0.001}{\alpha_{\rm eff}}\right)^{1/4}{\rm K},
\end{equation}
where 
\begin{equation}
\Omega_K=\sqrt{\frac{1}{R}\frac{\rd\Phi}{\rd R}}=\frac{c}{R_g}\sqrt{\frac{1}{r}}\frac{1}{r-2}.
\end{equation}

\section{Dimensionless Equations}
The variables in dimensionless forms are
\begin{gather}
\Omega_K=\frac{1}{r-2}\sqrt{\frac{1}{r}},~~~~~
\frac{\rd\ln\Omega_{\rm K}}{\rd r}=-\frac{1}{2r}-\frac{1}{r-2}.\\
W=\frac{\Sigma kT_0}{\bar\mu m_H}\frac{1}{c^2} + \frac{2HaT_0^4}{3}\frac{R_g^2}{\dot M_{\rm Edd}c}.\\
Q_{\rm rad}^-=\frac{8c(1-\beta)}{\kappa}\frac{W}{\Sigma H}
\end{gather}

For given $T$ and $\Omega$, we have 
\begin{equation}
W_g^\mu W^{\mu}
\end{equation}

\begin{thebibliography}{99}
\bibitem[Cao(2009)]{Cao2009} 
   Cao, X.\ 2009, \mnras, 394, 207 
\bibitem[Gong et al.(2012)]{Gong2012} 
   Gong, X.-L., Li, L.-X., \& Ma, R.-Y.\ 2012, \mnras, 420, 1415  
\bibitem[Coppi \& Blandford(1990)]{Coppi1990} 
   Coppi, P.~S., \& Blandford, R.~D.\ 1990, \mnras, 245, 453  
\bibitem[Kato et al.(2008)]{Kato2008}
   Kato, S., Fukue, J., \& Mineshige, S. 2008, Black Hole Accretion Disks—
   Towards a New Paradigm (Kyoto: Kyoto Univ. Press)   
\bibitem[Liu et al.(2002)]{Liu2002} 
   Liu, B.~F., Mineshige, S., \& Shibata, K.\ 2002, \apjl, 572, L173 
\bibitem[Merloni(2003)]{Merloni2003} 
   Merloni, A.\ 2003, \mnras, 341, 1051    
\bibitem[Narayan \& Yi(1995)]{Narayan1995} 
   Narayan, R., \& Yi, I.\ 1995, \apj, 452, 710    
\bibitem[Rybicki \& Lightman(1979)]{Rybicki1979} 
   Rybicki, G.~B., \& Lightman, A.~P.\ 1979, New York, Wiley-Interscience, 1979     
\bibitem[Svensson \& Zdziarski(1994)]{Svensson1994}
   Svensson, R., \& Zdziarski, A.~A.\ 1994, \apj, 436, 599  
\bibitem[Wang \& Netzer(2003)]{Wang2003} 
   Wang, J.-M., \& Netzer, H.\ 2003, \aap, 398, 927 
\bibitem[Wang \& Zhou(1999)]{Wang1999}
   Wang, J.-M., \& Zhou, Y.-Y.\ 1999, \apj, 516, 420    
\bibitem[Zycki et al.(1995)]{Zycki1995} 
   Zycki, P.~T., Collin-Souffrin, S., \& Czerny, B.\ 1995, \mnras, 277, 70 
\end{thebibliography}
\end{document}
